% ##############################################################################
\section{Policy Model}
\label{sec:policy_model}

The policy model has two layers. The first decides which modules exist for a tenant and role. The second decides which features and actions are available inside those modules.

\subsection{Modules and layout as data}
\label{sec:module_policy}

Module access is stored per tenant and role in DynamoDB. A policy record contains an enabled-module list and a \code{moduleLayout} object with structured UI fragments. The control plane reads this record during \code{/v2/me} materialization and merges tenant fragments with shipped defaults.

Tenant changes are therefore small and reviewable. Enabling DataMap for a pilot tenant or switching a sidebar mode is a data change. Adding a new module type still requires code, since a new module brings routes, components, handlers, and documentation. That is the right boundary: routine tenant variability belongs in data; new product surface belongs in code review.

Malformed configuration is rejected during materialization. Partial configuration is allowed only when it can be safely merged with defaults. A tenant can override one sidebar group without redefining the whole navigation model.

\subsection{Server-built abilities}
\label{sec:abilities}

Fine-grained permissions are expressed with CASL, an isomorphic JavaScript authorization library that serializes ability rules between server and client \cite{casl_docs}. In our deployment the ability model covers three actions, \code{view}, \code{manage}, and \code{use}, across 25+ subjects. Those subjects include modules, generic pages, feature permissions, sub-module subjects, and config-derived subjects such as \code{module:config:integration:enabled}.

The server builds the ability object from the materialized profile. The client receives packed rules and uses them only for UX decisions: hiding buttons, disabling menu entries, and suppressing unavailable panels. Backend APIs still enforce authorization independently. A browser that tampers with packed rules may fool itself; it does not gain backend access.

\subsection{No hidden second policy language}
\label{sec:no_second_policy}

The trap in many products is accidental policy language. A frontend branch like \code{if (tenant === "acme")} looks harmless. Ten of them become an undocumented entitlement system.

\uiPolicy{} draws a hard line. Frontend code may render from a contract. It may not invent access. The contract may be generated by RBAC, ABAC, a policy engine, a module registry, or a mix of those. What matters is that the rendered UI consumes the materialized decision instead of re-creating it.

\begin{table}[H]
  \centering
  \small
  \begin{tabularx}{\textwidth}{X r X}
    \toprule
    \textbf{Policy surface} & \textbf{Scale} & \textbf{Runtime effect} \\
    \midrule
    Modules in policy store & 7 & Drives navigation, dashboard cards, and module route access. \\
    CASL permission subjects & 25+ & Controls module, feature, page, sub-module, and config-derived access. \\
    \code{moduleLayout} fields & 150+ & Shapes sidebar mode, module defaults, settings tabs, feedback surfaces, and per-module fragments. \\
    Sidebar modes & 6 & Lets tenants vary navigation behavior without code forks. \\
    Tenant onboardings requiring deployment & 0 & Normal onboarding is a policy-store update plus identity assignment. \\
    \bottomrule
  \end{tabularx}
  \caption{Policy surface in the production deployment. Counts are taken from the deployed UI and control-plane configuration.}
  \label{tab:policy_surface}
\end{table}
